\documentclass[11pt]{article}
\usepackage[utf8]{inputenc}
\usepackage[T1]{fontenc}
\usepackage[italian]{babel}
\usepackage{geometry}
\geometry{a4paper, margin=1in}
\usepackage{xcolor}
\usepackage{listings}
\usepackage{hyperref}

\usepackage{tcolorbox}
\tcbuselibrary{breakable}
\begin{document}

\begin{tcolorbox}[colback=blue!5!white,colframe=blue!75!black,title=\textbf{DOMANDA}, breakable]
4) Sostituire ``127.0.0.1'' prima con ``localhost'' e poi con ``pippo'' e osservare cosa succede
\end{tcolorbox}

\begin{tcolorbox}[colback=green!5!white,colframe=green!75!black,title=\textbf{RISPOSTA}, breakable]
Ecco cosa succede se nel file \texttt{clientUDP.c} modifichiamo l'indirizzo di destinazione (attualmente ``127.0.0.1'') all'interno della chiamata a \texttt{UDPSend}:

\begin{itemize}
    \item \textbf{Sostituendo con ``localhost'':}
    Non si osserva alcun cambiamento nel comportamento del programma. Il client e il server continueranno a comunicare correttamente. Questo accade perché la funzione di rete utilizzata sotto il cofano dalla libreria (\texttt{gethostbyname()}) interroga il sistema di risoluzione dei nomi (DNS locale o file \texttt{/etc/hosts}). Il nome ``localhost'' viene risolto automaticamente nell'indirizzo \texttt{IP} di loopback 127.0.0.1. Di conseguenza, i pacchetti arrivano a destinazione allo stesso identico modo.
    
    \item \textbf{Sostituendo con ``pippo'':}
    Il client terminerà l'esecuzione con un errore (probabilmente in modo silenzioso o stampando un codice di errore, a seconda di come è gestita la libreria) prima di inviare il messaggio. Questo avviene perché la funzione \texttt{gethostbyname(``pippo'')} tenterà di risolvere l'hostname ``pippo'' in un indirizzo \texttt{IP}, ma fallirà (poiché ``pippo'' non è un nome noto nella rete locale, a meno di non averlo configurato manualmente nel file \texttt{/etc/hosts}). Andando a ispezionare la libreria \texttt{network.c}, si può notare che quando la risoluzione fallisce, la funzione \texttt{UDPSend} esegue una \texttt{exit(ERRHOST)} terminando brutalmente il processo. Il server, di conseguenza, rimarrà in perenne attesa del messaggio.
\end{itemize}
\end{tcolorbox}

\end{document}
